\section{Availability gap and data-center context}

The LD position cannot be set in isolation from the end-customer requirement.
Data-center loads are specified against the Uptime Institute tier framework, and
there is a structural gap between what a single HVDC link delivers and what those
frameworks demand.

\subsection{The tier requirement}
\begin{itemize}\setlength{\itemsep}{0.25em}
  \item The industry norm for large data centers is \textbf{Tier~III}
        ($\approx$99.982\,\% availability, $\approx$1.6\,h downtime/year).
  \item Tier~IV (full redundancy) is generally reserved for financial
        institutions and is not the target here.
  \item Banks and project financiers require this availability analysis as a
        condition of financing.
\end{itemize}

\subsection{The gap}
\begin{itemize}\setlength{\itemsep}{0.25em}
  \item A single HVDC link \textbf{cannot} meet the Tier~III requirement on its
        own.
  \item Even dual symmetrical-monopole configurations must include
        \textbf{tower/line failure} in the availability calculation --- a factor
        confirmed by the DC-line development team, whose assumed tower-failure
        rate was significantly higher than expected, and one that is easily
        overlooked (it is not always captured in single-contingency studies, but
        most ISOs include it in their evaluation of a double circuit).
\end{itemize}

\subsection{Tower failure and corridor sharing}
Where redundancy is provided by two circuits, tower failure becomes a key
consideration:
\begin{itemize}\setlength{\itemsep}{0.25em}
  \item \textbf{Tower spacing.} If both circuits share a structure, a single
        tower failure takes out both. The towers (or lines) must therefore be
        spaced far enough apart that a \emph{falling tower does not affect the
        adjacent circuit}.
  \item \textbf{Shared corridor.} Spacing alone is not sufficient: as long as
        both circuits remain in the \textbf{same corridor / right-of-way}, a
        \textbf{total outage remains possible from a common-cause event} such as
        wildfire or another natural disaster. A wildfire, for example, can take
        out a parallel circuit hundreds of metres --- even a mile --- away
        through smoke/pollution flashover; a tornado through a shared corridor
        has historically taken out both bipoles in the same right-of-way
        (Manitoba experience, which led to the third bipole being routed
        separately).
  \item \textbf{Implication.} In areas of significant wildfire or severe-weather
        risk, true redundancy may require \emph{separate corridors} or
        \emph{undergrounding} the circuits --- at a significant cable cost
        premium, though still potentially cheaper than local generation.
\end{itemize}

\subsection{Mitigation paths}
\begin{enumerate}\setlength{\itemsep}{0.2em}
  \item \textbf{Dual HVDC circuits}, each sized for full capacity (cf.\ the
        Johan Sverdrup~II approach).
  \item \textbf{On-site back-up generation} --- effective but costly and a
        challenge to the scheme's green credentials.
  \item \textbf{Negotiate the availability target down} with the data center ---
        feasibility depends on the counterparty and its financing.
\end{enumerate}

The chosen mitigation directly shapes the LD regime: the more the redundancy
strategy relies on the HVDC supplier's performance, the more weight the LDs must
carry.

\subsection{Manned vs.\ unmanned converter station}
Some operators argue that permanently manning a converter station improves
availability, since staff on site can respond faster to an event. In practice
this benefit is doubtful for a single project:
\begin{itemize}\setlength{\itemsep}{0.25em}
  \item \textbf{Motivation.} It is very difficult to keep qualified personnel
        motivated at a converter station: genuine incidents are rare and the
        role is largely routine maintenance, so skilled staff become bored and
        do not stay.
  \item \textbf{Skills.} Those willing to remain in such a role tend not to be
        the people who can provide a proper response in a real emergency; and
        once skilled staff leave, replacements have little of the requisite
        expertise --- so a permanently manned station adds cost without reliably
        adding capability.
  \item \textbf{Better approach.} Retain \emph{access} to skilled personnel with
        a plan to mobilise them quickly to the site, and/or an OEM service-level
        agreement covering maintenance, rather than permanently manning the
        station. Building an in-house maintenance team only becomes worthwhile
        across a portfolio of multiple DC projects.
\end{itemize}
